\documentclass[10pt,a4paper]{amsart}
\usepackage{amsmath,amssymb,amsthm,graphicx}
\usepackage{verbatim}
\newenvironment{reference}{\comment}{\endcomment}
%\newenvironment{reference}{}{}
\newenvironment{slogan}{\comment}{\endcomment}
\newenvironment{history}{\comment}{\endcomment}

% For commutative diagrams we use Xy-pic
\usepackage[all]{xy}

% We use 2cell for 2-commutative diagrams.
\xyoption{2cell}
% Theorem environments.
%
\theoremstyle{plain}
\newtheorem{theorem}[subsection]{Theorem}
\newtheorem{proposition}[subsection]{Proposition}
\newtheorem{lemma}[subsection]{Lemma}

\theoremstyle{definition}
\newtheorem{definition}[subsection]{Definition}
\newtheorem{example}[subsection]{Example}
\newtheorem{exercise}[subsection]{Exercise}
\newtheorem{situation}[subsection]{Situation}

\theoremstyle{remark}
\newtheorem{remark}[subsection]{Remark}
\newtheorem{remarks}[subsection]{Remarks}

\numberwithin{equation}{subsection}

% Macros
%
\def\lim{\mathop{\mathrm{lim}}\nolimits}
\def\colim{\mathop{\mathrm{colim}}\nolimits}
\def\Spec{\mathop{\mathrm{Spec}}}
\def\Hom{\mathop{\mathrm{Hom}}\nolimits}
\def\Ext{\mathop{\mathrm{Ext}}\nolimits}
\def\SheafHom{\mathop{\mathcal{H}\!\mathit{om}}\nolimits}
\def\SheafExt{\mathop{\mathcal{E}\!\mathit{xt}}\nolimits}
\def\Sch{\mathit{Sch}}
\def\Mor{\mathop{\mathrm{Mor}}\nolimits}
\def\Ob{\mathop{\mathrm{Ob}}\nolimits}
\def\Sh{\mathop{\mathit{Sh}}\nolimits}
\def\NL{\mathop{N\!L}\nolimits}
\def\CH{\mathop{\mathrm{CH}}\nolimits}
\def\proetale{{pro\text{-}\acute{e}tale}}
\def\etale{{\acute{e}tale}}
\def\QCoh{\mathit{QCoh}}
\def\Ker{\mathop{\mathrm{Ker}}}
\def\Im{\mathop{\mathrm{Im}}}
\def\Coker{\mathop{\mathrm{Coker}}}
\def\Coim{\mathop{\mathrm{Coim}}}

% Boxtimes
%
\DeclareMathSymbol{\boxtimes}{\mathbin}{AMSa}{"02}

%
% Macros for moduli stacks/spaces
%
\def\QCohstack{\mathcal{QC}\!\mathit{oh}}
\def\Cohstack{\mathcal{C}\!\mathit{oh}}
\def\Spacesstack{\mathcal{S}\!\mathit{paces}}
\def\Quotfunctor{\mathrm{Quot}}
\def\Hilbfunctor{\mathrm{Hilb}}
\def\Curvesstack{\mathcal{C}\!\mathit{urves}}
\def\Polarizedstack{\mathcal{P}\!\mathit{olarized}}
\def\Complexesstack{\mathcal{C}\!\mathit{omplexes}}
% \Pic is the operator that assigns to X its picard group, usage \Pic(X)
% \Picardstack_{X/B} denotes the Picard stack of X over B
% \Picardfunctor_{X/B} denotes the Picard functor of X over B
\def\Pic{\mathop{\mathrm{Pic}}\nolimits}
\def\Picardstack{\mathcal{P}\!\mathit{ic}}
\def\Picardfunctor{\mathrm{Pic}}
\def\Deformationcategory{\mathcal{D}\!\mathit{ef}}
\usepackage{cleveref}
\crefname{originaldefinition}{original definition}{original definitions}
\Crefname{originaldefinition}{Original definition}{Original definitions}
\makeatletter
\newlabel{bootstrap-definition-property-transformation}{{03XZ}{source}{}{sourcepointer-1}{}}
\newlabel{bootstrap-definition-property-transformation@cref}{{[originaldefinition][0][]03XZ}{[1][0][]source}{}{}{}}
\makeatother
\crefname{originalreference}{original reference}{original references}
\Crefname{originalreference}{Original reference}{Original references}
\makeatletter
\newlabel{definition-flat}{{00HB}{source}{}{sourcepointer-2}{}}
\newlabel{definition-flat@cref}{{[originalreference][0][]00HB}{[1][0][]source}{}{}{}}
\makeatother
\crefname{originallemma}{original lemma}{original lemmas}
\Crefname{originallemma}{Original lemma}{Original lemmas}
\makeatletter
\newlabel{flat-lemma-etale-flat-up-down}{{05B9}{source}{}{sourcepointer-3}{}}
\newlabel{flat-lemma-etale-flat-up-down@cref}{{[originallemma][0][]05B9}{[1][0][]source}{}{}{}}
\makeatother
\crefname{originallemma}{original lemma}{original lemmas}
\Crefname{originallemma}{Original lemma}{Original lemmas}
\makeatletter
\newlabel{formal-spaces-lemma-adic-homomorphism}{{0GBS}{source}{}{sourcepointer-4}{}}
\newlabel{formal-spaces-lemma-adic-homomorphism@cref}{{[originallemma][0][]0GBS}{[1][0][]source}{}{}{}}
\makeatother
\crefname{originallemma}{original lemma}{original lemmas}
\Crefname{originallemma}{Original lemma}{Original lemmas}
\makeatletter
\newlabel{formal-spaces-lemma-representable-local-property}{{0ANK}{source}{}{sourcepointer-5}{}}
\newlabel{formal-spaces-lemma-representable-local-property@cref}{{[originallemma][0][]0ANK}{[1][0][]source}{}{}{}}
\makeatother
\crefname{originallemma}{original lemma}{original lemmas}
\Crefname{originallemma}{Original lemma}{Original lemmas}
\makeatletter
\newlabel{more-algebra-lemma-limit-flat}{{0912}{source}{}{sourcepointer-6}{}}
\newlabel{more-algebra-lemma-limit-flat@cref}{{[originallemma][0][]0912}{[1][0][]source}{}{}{}}
\makeatother
\title{Representable flatness implies flatness for morphisms of formal algebraic spaces}
\author{The Stacks Project Authors}
\date{stacks-project:restricted.tex:local-snapshot-20261012; GNU FDL 1.2 or later}
\begin{document}
\maketitle
\noindent\textit{Editorial note.} All delivered mathematical text is the author's own, copied from the pinned Stacks Project chapter file restricted.tex; no statement or proof step was rewritten, reordered or removed.\par\medskip
\noindent\textit{Editorial note.} The counted proof is the author's complete argument: the source proves the lemma by a four-implication flatness chain on affine charts and closes with the limit flatness lemma; the argument is complete\par\medskip
\noindent\textit{Editorial note.} The label lemma-representable-flat names the counted statement, and the proof environment immediately following it in the source is delivered verbatim as its proof.\par\medskip
\noindent\textit{Editorial note.} The counted statement is carried by the author's own lemma environment in the chapter source.\par\medskip
\noindent\textit{Allowed standard prerequisites (editorial).}\begin{itemize}
\item morphisms of locally Noetherian formal algebraic spaces representable by algebraic spaces and flat in the transformation sense of the bootstrap chapter
\item flatness of morphisms of formal algebraic spaces defined via the taut maps of the WAdm\textasciicircum{}Noeth category on affine etale charts
\item etale morphisms are flat and flatness composes
\item the etale-flat up-down lemma for flatness after etale base change
\item flatness of a limit of ring maps from flatness of the reductions
\end{itemize}\medskip
\ifdefined\hypertarget\hypertarget{sourcepointer-1}{}\fi
\noindent\textit{Source reference.} Original definition 03XZ: sources/stacks-chow/tags.txt, source lines 5166--5166. Prerequisite pointer invoked at proof line(s) 3141; it names a standard fact of a sibling chapter that the counted argument uses, and every step of the counted proof built on it is retained verbatim in the delivered body.\par\medskip
\ifdefined\hypertarget\hypertarget{sourcepointer-2}{}\fi
\noindent\textit{Source reference.} Original reference 00HB: sources/stacks-restricted/pinned-original-source.tex, source lines 3023--3023. Prerequisite pointer invoked at proof line(s) 3142; it names a standard fact of the same chapter that the counted argument uses, and every step of the counted proof built on it is retained verbatim in the delivered body.\par\medskip
\ifdefined\hypertarget\hypertarget{sourcepointer-3}{}\fi
\noindent\textit{Source reference.} Original lemma 05B9: sources/stacks-chow/tags.txt, source lines 6877--6877. Prerequisite pointer invoked at proof line(s) 3194; it names a standard fact of a sibling chapter that the counted argument uses, and every step of the counted proof built on it is retained verbatim in the delivered body.\par\medskip
\ifdefined\hypertarget\hypertarget{sourcepointer-4}{}\fi
\noindent\textit{Source reference.} Original lemma 0GBS: sources/stacks-chow/tags.txt, source lines 20382--20382. Prerequisite pointer invoked at proof line(s) 3165; it names a standard fact of a sibling chapter that the counted argument uses, and every step of the counted proof built on it is retained verbatim in the delivered body.\par\medskip
\ifdefined\hypertarget\hypertarget{sourcepointer-5}{}\fi
\noindent\textit{Source reference.} Original lemma 0ANK: sources/stacks-chow/tags.txt, source lines 13443--13443. Prerequisite pointer invoked at proof line(s) 3163; it names a standard fact of a sibling chapter that the counted argument uses, and every step of the counted proof built on it is retained verbatim in the delivered body.\par\medskip
\ifdefined\hypertarget\hypertarget{sourcepointer-6}{}\fi
\noindent\textit{Source reference.} Original lemma 0912: sources/stacks-chow/tags.txt, source lines 11430--11430. Prerequisite pointer invoked at proof line(s) 3205; it names a standard fact of a sibling chapter that the counted argument uses, and every step of the counted proof built on it is retained verbatim in the delivered body.\par\medskip
\begin{lemma}
\label{lemma-representable-flat}
Let $S$ be a scheme. Let $f : X \to Y$ be a morphisms of
locally Noetherian formal algebraic spaces over $S$.
If $f$ is representable by algebraic spaces and
flat in the sense of Bootstrap, Definition
\ref{bootstrap-definition-property-transformation},
then $f$ is flat in the sense of Definition \ref{definition-flat}.
\end{lemma}

\begin{proof}
This is a sanity check whose proof should be trivial but isn't quite.
We urge the reader to skip the proof.
Assume $f$ is representable by algebraic spaces and
flat in the sense of Bootstrap, Definition
\ref{bootstrap-definition-property-transformation}.
Consider a commutative diagram
$$
\xymatrix{
U \ar[d] \ar[r] & V \ar[d] \\
X \ar[r] & Y
}
$$
with $U$ and $V$ affine formal algebraic spaces, $U \to X$ and $V \to Y$
representable by algebraic spaces and \'etale.
Then the morphism $U \to V$
corresponds to a taut map $B \to A$ of $\textit{WAdm}^{Noeth}$ by
Formal Spaces, Lemma
\ref{formal-spaces-lemma-representable-local-property}.
Observe that this means $B \to A$ is adic (Formal Spaces, Lemma
\ref{formal-spaces-lemma-adic-homomorphism})
and in particular for any ideal of definition $J \subset B$
the topology on $A$ is the $J$-adic topology and the diagrams
$$
\xymatrix{
\Spec(A/J^nA) \ar[r] \ar[d] & \Spec(B/J^n) \ar[d] \\
U \ar[r] & V
}
$$
are cartesian.

\medskip\noindent
Let $T \to V$ is a morphism where $T$ is a scheme. Then
\begin{align*}
X \times_Y T \to T\text{ is flat}
& \Rightarrow
U \times_Y T \to T\text{ is flat} \\
& \Rightarrow
U \times_V V \times_Y T \to T\text{ is flat} \\
& \Rightarrow
U \times_V V \times_Y T \to V \times_Y T\text{ is flat} \\
& \Rightarrow
U \times_V T \to T\text{ is flat}
\end{align*}
The first statement is the assumption on $f$.
The first implication because $U \to X$ is \'etale and hence flat
and compositions of flat morphisms of algebraic spaces are flat.
The second impliciation because $U \times_Y T = U \times_V V \times_Y T$.
The third implication by
More on Flatness, Lemma \ref{flat-lemma-etale-flat-up-down}.
The fourth implication because we can pullback by the morphism
$T \to V \times_Y T$.
We conclude that $U \to V$ is flat in the sense of
Bootstrap, Definition
\ref{bootstrap-definition-property-transformation}.
In terms of the continuous ring map $B \to A$
this means the ring maps $B/J^n \to A/J^nA$ are flat (see diagram above).

\medskip\noindent
Finally, we can conclude that $B \to A$ is flat for example
by More on Algebra, Lemma \ref{more-algebra-lemma-limit-flat}.
\end{proof}

\end{document}
